\documentclass[11pt]{article}
\usepackage[T1]{fontenc}
\usepackage{lmodern}
\usepackage[margin=1in]{geometry}
\usepackage{amsmath,amssymb,amsthm,mathtools}
\usepackage{microtype,booktabs,array,enumitem}
\usepackage[hidelinks]{hyperref}
\hypersetup{pdftitle={IE-28: Recovered Analysis and a Three-Stage Existence Proof},
 pdfauthor={ChatGPT},pdfsubject={Partial resolution: two and three stages; general case not proved},
 pdfkeywords={collocation, diagonal preconditioning, nilpotence, IE-28}}
\usepackage{fancyhdr}
\setlength{\headheight}{14pt}
\pagestyle{fancy}
\fancyhf{}
\lhead{\small IE-28: recovered analysis}
\rhead{\small Two- and three-stage existence}
\cfoot{\thepage}
\setlength{\parskip}{4pt}
\setlength{\parindent}{0pt}
\setlist{nosep,leftmargin=*}
\newtheorem{theorem}{Theorem}[section]
\newtheorem{lemma}[theorem]{Lemma}
\newtheorem{proposition}[theorem]{Proposition}
\newtheorem{corollary}[theorem]{Corollary}
\theoremstyle{remark}
\newtheorem{remark}[theorem]{Remark}
\DeclareMathOperator{\diag}{diag}
\DeclareMathOperator{\tr}{tr}
\DeclareMathOperator{\spec}{spec}
\newcommand{\R}{\mathbb R}
\newcommand{\one}{\mathbf 1}
\title{\textbf{IE-28: Recovered Analysis and a\\Three-Stage Existence Proof}\\[6pt]
\large Positive diagonal preconditioning with a nilpotent stiff limit}
\author{Prepared by ChatGPT}
\date{12 September 2026}

\begin{document}
\maketitle
\thispagestyle{empty}

\begin{abstract}
For the collocation integration matrix on arbitrary distinct positive nodes, we
reconstruct the inverse-matrix and eigenpolynomial reductions from an interrupted
attempt. We give complete, self-contained existence proofs for two and three
stages. In the three-stage case a determinant-normalized family of cubic
polynomials forces an eigenvalue at one, and two exact trace identities yield a
sign change. The resulting diagonal is positive, increasingly ordered, and obeys
$c_i/3<d_i<c_i$. The general existence statement for arbitrary stage count is
\emph{not proved in this document}. Higher-stage calculations are reported only
as numerical candidates. The archive contains executable symbolic checks,
high-precision numerical code, and rerun results; it is not an all-stage solution
or a peer-reviewed claim of priority.
\end{abstract}

\section*{Scope and recovery status}
The original hyperlink points to entry \textbf{IE-28} in
\emph{OpenProblemsInNLA}~\cite{entry}. This document follows that destination,
not the inconsistent visible link label in the request. The earlier working
files were absent when recovery began. Code and mathematical developments were
reconstructed from the surviving conversation record; all result files in this
archive were generated again. The three-stage sign-change argument is completed
here rather than represented as an already-finished recovered proof.

\begin{center}
\begin{tabular}{@{}p{0.19\textwidth}p{0.72\textwidth}@{}}
\toprule
Scope & What this archive establishes \\
\midrule
$s=2$ & An explicit positive diagonal for every allowed node pair. \\
$s=3$ & A proof for every allowed node triple, with a bracketed scalar construction. \\
General $s$ & Exact algebraic reductions and a positive eigenpolynomial ansatz. \\
$s\geq4$ & Selected high-precision candidates; no universal existence proof or counterexample. \\
\bottomrule
\end{tabular}
\end{center}

\textbf{The missing part of the original request is precise:} a proof of existence
for every stage count $s\geq4$ and every admissible node set is not supplied.
The proofs below do not depend on the numerical experiments.

\newpage
\tableofcontents
\medskip

\section{Target and algebraic equivalences}
Let $s\geq2$ and $0<c_1<\cdots<c_s\leq1$. Define
\begin{equation}\label{eq:problem}
 \ell_j(t)=\prod_{k\ne j}\frac{t-c_k}{c_j-c_k},\qquad
 A_{ij}=\int_0^{c_i}\ell_j(t)\,dt.
\end{equation}
The target is a real diagonal matrix $D=\diag(d_1,\ldots,d_s)$, with each $d_i>0$,
such that
\begin{equation}\label{eq:target}
 N:=I-D^{-1}A,\qquad N^s=0.
\end{equation}
The positive-diagonal zero-spectral-radius conjecture appears in van der Houwen
and de Swart~\cite[\S3.2.1, p.~46]{houwen}. The 2025 MIN-SR-S treatment explicitly
leaves existence at arbitrary stage count unguaranteed~\cite[\S2.2.3, p.~A439]{sdc}.
These are context references, not substitutes for the arguments below.

Put
\[
 C=\diag(c_1,\ldots,c_s),\quad V_{ik}=c_i^{k-1},\quad
 R=\diag(1,1/2,\ldots,1/s).
\]
Integrating each monomial of degree at most $s-1$ gives
\begin{equation}\label{eq:vandermonde}
 A=CVRV^{-1},\qquad \det A=\frac{\prod_{i=1}^s c_i}{s!}>0.
\end{equation}
In particular $A$ is nonsingular. Scaling every node by $\gamma>0$ scales $A$
and a corresponding $D$ by $\gamma$, leaving $D^{-1}A$ unchanged. The upper
bound $c_s\leq1$ therefore only fixes a scale.

Let $B=A^{-1}D=(D^{-1}A)^{-1}$. Then \eqref{eq:target} is equivalent to
\begin{equation}\label{eq:char}
 \det(zI-B)=(z-1)^s.
\end{equation}
Indeed, if $B=I+J$ with $J^s=0$, then
$N=I-B^{-1}=B^{-1}J$, and the factors commute; thus $N^s=0$.
Conversely, $N^s=0$ gives
$B=(I-N)^{-1}=I+N+\cdots+N^{s-1}$, so $B-I$ is nilpotent.
Cayley--Hamilton converts between this statement and \eqref{eq:char}.
Every solution necessarily satisfies
\begin{equation}\label{eq:prod}
 \prod_{i=1}^s d_i=\det A=\frac{\prod_{i=1}^s c_i}{s!}.
\end{equation}
The target is also equivalent to $\det((1-t)D+tA)=\det D$ for all real $t$: after factoring
$D$, the left side becomes $\det D\det(I-tN)$.

\section{An inverse formula and an eigenpolynomial construction}
\subsection{A diagonal similarity with simple off-diagonal entries}
Let
\[
 w_i=\frac{1}{\prod_{k\ne i}(c_i-c_k)},\qquad
 S=\diag(c_1/w_1,\ldots,c_s/w_s).
\]
The entries of $S$ need not all have the same sign, but none vanishes.
Define
\begin{equation}\label{eq:M}
 M_{ij}=\begin{cases}
 \displaystyle\frac{1}{c_i-c_j},&i\ne j,\\[6pt]
 \displaystyle h_i:=\frac1{c_i}+\sum_{k\ne i}\frac1{c_i-c_k},&i=j.
 \end{cases}
\end{equation}
\begin{lemma}\label{lem:inverse}
$A^{-1}=SMS^{-1}$. Consequently $L:=MD$ is similar to $B=A^{-1}D$,
and $\det M=s!/\prod_i c_i$.
\end{lemma}
\begin{proof}
The polynomials $g_j(t)=t\ell_j(t)/c_j$ have degree at most $s$, vanish at zero,
and satisfy $g_j(c_i)=\delta_{ij}$. Differentiation of the unique interpolant in
$\{p:\deg p\leq s,\ p(0)=0\}$ is the inverse of the integration map $A$.
Thus $(A^{-1})_{ij}=g_j'(c_i)$. At $i\ne j$ this equals
\[
 \frac{c_i}{c_j}\frac{w_j}{w_i}\frac1{c_i-c_j}
 =\frac{S_{ii}}{S_{jj}}M_{ij}.
\]
At $i=j$ it equals $1/c_i+\ell_i'(c_i)=h_i$. Since $S$ and $D$ commute,
$A^{-1}D=S(MD)S^{-1}$. The determinant assertion follows from
\eqref{eq:vandermonde}.
\end{proof}
The trace has the useful symmetric form
\begin{equation}\label{eq:tracepair}
 \tr(MD)=\sum_i\frac{d_i}{c_i}
       +\sum_{i<j}\frac{d_i-d_j}{c_i-c_j}.
\end{equation}
This follows by grouping the two terms for each pair of nodes in
$\sum_i h_i d_i$; no sign assumption on the individual $h_i$ is needed.

\subsection{Forcing one eigenvalue exactly}
\begin{lemma}[Eigenpolynomial]\label{lem:poly}
Suppose $p$ is a nonconstant polynomial of degree at most $s$, $p(0)=0$, and
$p'(c_i)\ne0$ for every $i$. Set $d_i=p(c_i)/p'(c_i)$. Then $A^{-1}D$, and
therefore $MD$, has eigenvalue one.
\end{lemma}
\begin{proof}
Write $q_i=p'(c_i)$. Interpolation is exact for $p'$, since its degree is at
most $s-1$. Therefore $(Aq)_i=\int_0^{c_i}p'(t)\,dt=p(c_i)=d_iq_i$.
The vector $q$ is nonzero, and $A^{-1}Dq=q$.
\end{proof}
A particularly useful positive family is
\begin{equation}\label{eq:ansatz}
 p(t)=t\prod_{j=1}^{s-1}(t+\alpha_j),\quad \alpha_j>0,\qquad
 d_i=\left(\frac1{c_i}+\sum_{j=1}^{s-1}\frac1{c_i+\alpha_j}\right)^{-1}.
\end{equation}
For this family
\begin{equation}\label{eq:bounds}
 \frac{c_i}{s}<d_i<c_i,\qquad d_1<d_2<\cdots<d_s.
\end{equation}
For the first assertion, $0<1/(c_i+\alpha_j)<1/c_i$ for every $j$. For the second, the positive function
$d(t)=(1/t+\sum_j1/(t+\alpha_j))^{-1}$ has derivative
\[
 d'(t)=d(t)^2\left(\frac1{t^2}+\sum_j\frac1{(t+\alpha_j)^2}\right)>0.
\]
These are properties of the ansatz, not necessary conditions imposed on every
possible solution of IE-28.

\section{Complete two-stage construction}
\begin{theorem}\label{thm:two}
For any $0<a<b$, let $r=\sqrt{2ab}$ and
\begin{equation}\label{eq:two}
 d_1=\frac{a(a+r)}{2a+r},\qquad d_2=\frac{b(b+r)}{2b+r}.
\end{equation}
For the two-node collocation matrix, $(I-D^{-1}A)^2=0$. Moreover,
$a/2<d_1<a$, $b/2<d_2<b$, and $d_1<d_2$.
\end{theorem}
\begin{proof}
Use $p(t)=t(t+r)$ in Lemma~\ref{lem:poly}. Thus one eigenvalue of $MD$ is one.
Also
\[
 \det(MD)=\frac{2(a+r)(b+r)}{(2a+r)(2b+r)}=1,
\]
because the difference between numerator and denominator is $r^2-2ab=0$.
The other eigenvalue is therefore one as well. Equation~\eqref{eq:char} proves
the conclusion. The inequalities follow from \eqref{eq:bounds}.
\end{proof}
The first entry in \eqref{eq:two} is algebraically equivalent to the elementary
formula displayed in the repository entry; the present derivation fits the
same eigenpolynomial framework as the next section.

\section{Complete three-stage existence proof}
\begin{theorem}\label{thm:three}
For every $0<c_1<c_2<c_3$, there exist $r>0$ and $0<\theta<1$ such that
\begin{equation}\label{eq:threeD}
 d_i=\left(\frac1{c_i}+\frac1{c_i+r}
                  +\frac{\theta}{\theta c_i+r}\right)^{-1}
\end{equation}
satisfies $(I-D^{-1}A)^3=0$. In addition,
\[
 0<d_1<d_2<d_3,\qquad c_i/3<d_i<c_i\quad(i=1,2,3).
\]
\end{theorem}
We prove the theorem by a determinant-normalized deformation with a scalar
trace sign change. Every step works for arbitrary positive distinct real nodes.

\subsection{A compact shape parameter and determinant normalization}
For $r>0$ and $0\leq\theta\leq1$, set
\begin{equation}\label{eq:p3}
 p_{r,\theta}(t)=t(t+r)(\theta t+r).
\end{equation}
This polynomial has degree three for $\theta>0$ and degree two at $\theta=0$.
In both cases its logarithmic derivative gives \eqref{eq:threeD}, and
Lemma~\ref{lem:poly} forces one eigenvalue of $L(r,\theta)=M D(r,\theta)$ to
be one. Define
\begin{equation}\label{eq:G}
 G(r,\theta)=\det L(r,\theta)
 =6\prod_{i=1}^3\left(1+\frac{c_i}{c_i+r}
                 +\frac{\theta c_i}{\theta c_i+r}\right)^{-1}.
\end{equation}
For fixed $\theta$, this is continuous and strictly increasing in $r>0$.
The parenthesized denominators strictly decrease with $r$. Its endpoint
values are
\[
 \lim_{r\downarrow0}G(r,0)=\frac68<1,\qquad
 \lim_{r\downarrow0}G(r,\theta)=\frac6{27}<1\ (\theta>0),\qquad
 \lim_{r\to\infty}G(r,\theta)=6>1.
\]
Hence there is a unique positive $r=r(\theta)$ with $G(r(\theta),\theta)=1$.

This root depends continuously on $\theta\in[0,1]$. One elementary justification
is as follows. For fixed $r>0$, the quantity
$\theta c_i/(\theta c_i+r)$ strictly increases with $\theta$, so $G$ decreases
and $r(\theta)$ increases. Thus all the roots lie in the compact positive
interval $[r(0),r(1)]$. If $\theta_n\to\theta$, every convergent subsequence of
$r(\theta_n)$ must, by continuity of $G$, converge to its unique root at
$\theta$. This proves continuity of the whole root function.

There is also a convenient computational upper bound. At $r=4c_3$ each factor
inside parentheses in \eqref{eq:G} is at most $7/5$, since
$\theta c/(\theta c+r)\leq c/(c+r)\leq1/5$.
Thus $G(4c_3,\theta)\geq6(5/7)^3>1$, providing a uniform finite bracket.

\subsection{Two exact endpoint trace identities}
Let
\[
 \sigma_1=c_1+c_2+c_3,\quad
 \sigma_2=c_1c_2+c_1c_3+c_2c_3,\quad \sigma_3=c_1c_2c_3.
\]
The following identities are valid for every $r>0$. They can be obtained by
substitution into \eqref{eq:tracepair} and taking a common denominator; the
archive also verifies them with exact symbolic arithmetic.

At $\theta=0$ the entries are $d_i=c_i(c_i+r)/(2c_i+r)$, and
\begin{equation}\label{eq:positive}
 \sum_{i=1}^3 h_i\frac{c_i(c_i+r)}{2c_i+r}-3
 =\frac{r(3r^2+3\sigma_1r+2\sigma_2)}{\prod_i(2c_i+r)}>0.
\end{equation}
At $\theta=1$ the entries are $d_i=c_i(c_i+r)/(3c_i+r)$. Put
\begin{align}
 P_3(r)&=\prod_i(3c_i+r),\notag\\
 F_3(r)&=6\prod_i(c_i+r)-P_3(r)\notag\\
       &=5r^3+3\sigma_1r^2-3\sigma_2r-21\sigma_3.
\end{align}
Then
\begin{equation}\label{eq:negative}
 P_3(r)\left(\sum_{i=1}^3 h_i\frac{c_i(c_i+r)}{3c_i+r}-3\right)
       =F_3(r)-2(r^3+3\sigma_3).
\end{equation}
The determinant normalization at $r=r(1)$ is exactly $F_3(r(1))=0$.
The right-hand side of \eqref{eq:negative} is therefore strictly negative there.

\subsection{A sign change completes the spectrum}
The continuous real-valued function
\[
 T(\theta)=\tr L(r(\theta),\theta)-3
\]
has $T(0)>0$ by \eqref{eq:positive} and $T(1)<0$ by
\eqref{eq:negative}. The intermediate value theorem supplies a
$\theta_*\in(0,1)$ with $T(\theta_*)=0$. No monotonicity or uniqueness of the
trace zero is needed. Set $r_*=r(\theta_*)$ and use \eqref{eq:threeD}.
The resulting $3\times3$ matrix $L$ has
\[
 1\in\spec(L),\qquad \tr L=3,\qquad \det L=1.
\]
Writing its characteristic polynomial as $z^3-3z^2+e_2z-1$, evaluation at $z=1$
gives $e_2=3$. Thus it is exactly $(z-1)^3$. Lemma~\ref{lem:inverse} and
\eqref{eq:char} imply $(I-D^{-1}A)^3=0$.
Finally, $p_{r_*,\theta_*}$ is a positive scalar multiple of
$t(t+r_*)(t+r_*/\theta_*)$, with two positive, distinct negative-root magnitudes.
The inequalities and ordering follow from \eqref{eq:bounds}.
This proves Theorem~\ref{thm:three}. \hfill$\square$

\section{A reproducible construction and its numerical checks}
\subsection{Nested one-dimensional root finding}
The proof gives a direct algorithm. For each $\theta\in[0,1]$, solve the
strictly monotone equation $G(r,\theta)=1$ on $(0,4c_3)$; then bisect a sign
change of $T(\theta)$ on $[0,1]$. Only the inner root is asserted to be unique.
The implementation normalizes $c_3$ to one, uses safeguarded Newton steps for
the inner root, and uses bisection for the outer root. It raises an exception
rather than silently accepting a failed residual or lost endpoint sign.

For the inner Newton iteration let
$u_i=1+c_i/(c_i+r)+\theta c_i/(\theta c_i+r)$. The derivative used is
\[
 \frac{\partial}{\partial r}\log G(r,\theta)
 =\sum_i\frac{c_i/(c_i+r)^2+\theta c_i/(\theta c_i+r)^2}{u_i}>0.
\]
This confirms that the numerical step corresponds to the monotone equation
proved above, rather than an unrelated optimization objective.

\subsection{A trace evaluation without small node differences}
Formula \eqref{eq:tracepair} removes many cancelling terms, but a direct
subtraction $d_i-d_j$ can still lose digits. Write
\[
 f(t)=\frac{t(at^2+bt+g)}{3at^2+2bt+g},\qquad
 (a,b,g)=(\theta,(1+\theta)r,r^2).
\]
For $x\ne y$,
\begin{equation}\label{eq:divdiff}
 \frac{f(x)-f(y)}{x-y}
  =\frac{H(x,y)}{(3ax^2+2bx+g)(3ay^2+2by+g)},
\end{equation}
where
\begin{align*}
 H(x,y)={}&3a^2x^2y^2+2abxy(x+y)+ag(x^2+y^2)\\
          &+2(b^2-ag)xy+bg(x+y)+g^2.
\end{align*}
Here $a\geq0$, $b,g>0$, and
$b^2-ag=(1+\theta+\theta^2)r^2>0$. Hence the trace calculation can use a sum
of positive terms, with no division by $c_i-c_j$. This exact identity is
independently checked in \texttt{scripts/check\_symbolic.py}.

\subsection{Representative values and diagnostics}
For $(c_1,c_2,c_3)=(0.1,0.5,1)$, the constructed root is approximately
\begin{align*}
 \theta_*&=0.0575684366179453038651271952868,\\
 r_*&=0.147747942077272910567405294582,
\end{align*}
and
\[
 (d_1,d_2,d_3)\approx
 (0.0693895714943835768,\ 0.2584034251022362534,\ 0.4647573668722332063).
\]
The determinant-normalized endpoint trace differences are approximately
$0.1664848539500202$ and $-0.1003963765322660$. These displayed decimals are
rounded approximations, not the exact diagonal asserted by the theorem.

The checked high-precision run used 200-digit arithmetic for diagnostics. Define
$e_k(L)$ by $\det(I+tL)=\sum_{k=0}^s e_k(L)t^k$, and let
\begin{equation}\label{eq:err}
 E_{\rm char}=\max_{1\leq k\leq s}
       \left|\frac{e_k(L)}{\binom sk}-1\right|,\qquad
 E_N=\frac{\|N^s\|_\infty}{\max(1,\|N\|_\infty^s)}.
\end{equation}
The JSON results retain additional digits, the diagonal entries, and absolute
power norms. Selected measured characteristic-coefficient residuals are:
\begin{center}
\begin{tabular}{@{}lll@{}}
\toprule
Stages & Nodes / case & $E_{\rm char}$ \\
\midrule
3 & $(0.1,0.5,1)$ & $1.89\times10^{-116}$ (upper rounded) \\
3 & $(1-2\cdot10^{-12},1-10^{-12},1)$ & $2.55\times10^{-117}$ (upper rounded) \\
4 & $(0.1,0.3,0.6,1)$, numerical candidate & $4.94\times10^{-131}$ (upper rounded) \\
5 & $c_i=i/5$, numerical candidate & $1.74\times10^{-119}$ (upper rounded) \\
6 & $c_i=i/6$, numerical candidate & $2.58\times10^{-114}$ (upper rounded) \\
7 & $c_i=i/7$, numerical candidate & $1.41\times10^{-119}$ (upper rounded) \\
\bottomrule
\end{tabular}
\end{center}
A separate regression grid of all 66 triples $(a/13,b/13,1)$ with
$1\leq a<b\leq12$ passed, with maximum measured $E_{\rm char}<3.31\times10^{-51}$.
The seven individually recorded triple tests include extreme clusters and mixed
scales. These finite tests check the implementation; they are not the reason
the theorem holds for every real node triple.

Small computed residuals alone do not certify the existence of an exact nearby
root for higher stages. No interval arithmetic, interval Newton test, or other
root-existence certificate is supplied for those candidates. Their numerical
success is not promoted to an all-node theorem.

Direct spectral-radius estimates are especially easy to misinterpret here.
For example, a nilpotent Jordan block with an $\varepsilon$ entry added in its
lower-left corner has characteristic polynomial $z^s-\varepsilon$ and spectral
radius $|\varepsilon|^{1/s}$. Thus tiny coefficient perturbations can yield
visibly nonzero eigenvalues near a multiple nilpotent eigenvalue. The code uses
coefficient and power residuals and labels them as diagnostics, not proofs.

\section{The exact remaining higher-stage gap}
\subsection{Reduced equations for the negative-root ansatz}
For $p$ in \eqref{eq:ansatz}, one eigenvalue is already one. The remaining
conditions can be imposed as
\begin{equation}\label{eq:general}
 e_k(MD)=\binom sk\quad(1\leq k\leq s-2),\qquad
 s!\prod_i\frac{d_i}{c_i}=1.
\end{equation}
These are $s-1$ equations for the $s-1$ positive numbers $\alpha_j$.
They suffice: evaluating the characteristic polynomial at one determines the
only missing coefficient, $e_{s-1}=s$. In particular no coefficient constraint
has merely been omitted from the three-stage proof.

The reconstructed high-precision search writes $\alpha_j=\exp(x_j)$ and applies
a damped Newton method to \eqref{eq:general}. Its last residual is the logarithm
of the determinant condition. A positive numerical solution of these reduced
equations is a useful candidate; the following uniform assertion is \emph{not}
established here:
\[
 \text{For every admissible node set, \eqref{eq:general} has a solution in }
 (0,\infty)^{s-1}.
\]
Even proving this stronger negative-root-ansatz assertion would require an
argument absent from the recovered work. It has not been shown that every
positive nilpotent scaling must lie in this ansatz, either. A failure of this
particular search or ansatz would therefore not itself refute IE-28.

\subsection{Why the three-stage endpoint argument is not an all-stage proof}
The three scalar invariants used above do not suffice in general higher dimension.
For example, the real matrix
\[
 B_4=\diag\!\left(1,2,
  \begin{bmatrix}1/2&-1/2\\1/2&1/2\end{bmatrix}\right)
\]
has trace four, determinant one, and eigenvalue one, but
\[
 \det(zI-B_4)=(z-1)(z-2)(z^2-z+1/2)\ne(z-1)^4.
\]
This is an illustration of insufficient invariants, \textbf{not} a counterexample
to the collocation conjecture: no claim is made that $B_4$ is of the required
form $A^{-1}D$. At $s\geq4$, simultaneous control of further coefficients is
essential. Dimension counting, parameter positivity, and successful numerical
continuation do not by themselves provide it.

\subsection{Different limits must not be conflated}
For the diagonal-preconditioned iteration, the stiff limit is
$I-D^{-1}A$, whereas the leading nonstiff matrix is $A-D$. The nilpotence theorem
for the latter and a product formula involving different diagonals on successive
sweeps in the 2025 paper do not establish nilpotence for one fixed stiff-limit
diagonal~\cite[\S\S2.2.2--2.2.4]{sdc}. Likewise, a triangular factorization does
not meet the diagonal restriction. No such substitution is used in this note.

\section{Reproduction, provenance, and review boundary}
The recommended entry point is \texttt{README.md}. In a Python environment with
the supplied dependencies, run from the archive's root directory:
\begin{verbatim}
python scripts/check_symbolic.py
python scripts/run_examples.py
\end{verbatim}
The first script verifies the universal rational endpoint identities, the
cancellation-free divided difference, the two-stage reduction, and the cubic
completion identity. It additionally checks inverse similarity, determinant,
and eigenpolynomial formulas on exact rational examples at stages two through
six, with direct integration checks through stage four. The universal proofs of
the structural formulas are in Section~2; the finite rational tests are not
being presented as universal proofs of those formulas.

The second script reruns two explicit two-stage examples, seven three-stage
examples, six higher-stage numerical candidates, the 66-triple regression grid,
and input validation. It compares the three-stage answer with an independent
Newton refinement of the ansatz equations. It writes JSON and text reports in
\texttt{results/}. Original binary result files from the interrupted session
were not recovered and are not claimed to have been recovered.

The \texttt{experiments/} folder retains selected reconstructed exploratory
scripts. Their tolerances, unvalidated path tracking, and incomplete real-root
counting must not be read as certificates. They are separated from the proved
construction and its verification scripts. The proof is intended to be
self-contained and checkable; no claim of independent peer review or literature
priority is made.

\medskip
\textbf{Conclusion.} The recovery yields a complete constructive proof for two
and three stages, exact reductions for all stages, and reproducible computational
work. It does not finish the originally requested all-stage IE-28 assertion.

\begin{thebibliography}{9}
\bibitem{entry}
\emph{OpenProblemsInNLA}, IE-28, ``Positive diagonal preconditioning with a
nilpotent stiff limit,'' repository entry, accessed 12 September 2026.
\href{https://github.com/ajt60gaibb/OpenProblemsInNLA/blob/main/linear-systems-and-elimination/IE-28/README.md}{GitHub problem statement}.

\bibitem{houwen}
P.~J. van der Houwen and J.~J.~B. de Swart,
``Triangularly implicit iteration methods for ODE-IVP solvers,''
\emph{SIAM Journal on Scientific Computing} \textbf{18}(1) (1997), 41--55.
\href{https://doi.org/10.1137/S1064827595287456}{doi:10.1137/S1064827595287456}.
See \S3.2.1, p.~46, and \S3.2.2.
\href{https://ir.cwi.nl/pub/2191/2191D.pdf}{CWI full text}.

\bibitem{sdc}
G.~\v{C}aklovi\'{c}, T.~Lunet, S.~G\"otschel, and D.~Ruprecht,
``Improving efficiency of parallel across the method spectral deferred
corrections,'' \emph{SIAM Journal on Scientific Computing}
\textbf{47}(1) (2025), A430--A453.
\href{https://doi.org/10.1137/24M1649800}{doi:10.1137/24M1649800}.
See \S2.2.3, especially p.~A439.
\href{https://d-nb.info/1363153935/34}{Open full text}.
\end{thebibliography}
\end{document}
